% #############################################################################
% This is Chapter 1
% !TEX root = ../main.tex
% #############################################################################
% Change the Name of the Chapter i the following line
\fancychapter{Introduction}
% The following line allows to ref this chapter
\label{chap:intro}

Machine learning has become an important approach to network intrusion detection. 
Signature-based detectors rely on known attack patterns, which limits their ability to detect new attacks or variants that do not match existing signatures.
They can also be evaded by modifying traffic to avoid these patterns while preserving attack functionality.
Learned models can generalize beyond known signatures, but introduce a different attack surface: attackers can deliberately manipulate traffic features to exploit the models' learned decision boundaries.
This thesis addresses such adversarial evasion through changes to packet timing and size that preserve protocol compliance and attack functionality, and investigates how to defend against it.
It proposes DEPT
(Defense via Ensemble and Purification Transformer), a streaming defense that detects
adversarially manipulated network traffic and aims to reverse the manipulation before it reaches
the detector it was meant to evade.

% #############################################################################
\section{Context and Motivation}\label{sec:intro-context}
A \ac{NIDS} monitors network traffic to identify malicious activity.
Traditionally, detection has relied on matching traffic against signatures
of known threats, which can accurately identify known attacks but may miss new attacks or variants that do not match existing signatures.
\ac{ML} offers a different approach: instead of explicitly specifying attack patterns, a model learns the statistical behaviour that separates malicious from benign
traffic, and can therefore generalize to variants and previously unseen attacks~\cite{boutaba2018comprehensive,wang2021machine}. Learning-based detection is a major research direction in the field and an increasingly common deployment
choice~\cite{xu2025deep,anis2025survey}.

This shift introduced a new attack surface. A learned model can be fooled: a small,
deliberate perturbation of its input may cause a misclassification.
An input crafted to induce such an error is known as
\emph{adversarial example}~\cite{szegedy2014intriguing,goodfellow2014explaining}. For a
\ac{NIDS}, adversarial examples can enable \emph{evasion}: the attacker reshapes
ongoing attack traffic just enough for the detector to label it benign, while preserving the attack's functionality~\cite{he2023adversarial}. This threat is not hypothetical. Prior work has demonstrated that black-box attacks can evade deep-learning-based intrusion detectors by manipulating traffic features~\cite{zhang2020tiki}. \Cref{fig:intro-evasion}
illustrates the threat addressed in this thesis.

\begin{figure}[htbp]
    \centering
    \includesvg[width=0.98\textwidth]{./Images/intro_evasion}
    \caption[The adversarial evasion problem.]{The problem this thesis addresses. The original attack traffic is flagged by the ML-based
    NIDS that guards the network (top path). A black-box generator, TANTGen~\cite{pissarra2026tantgen} in this work,
    reshapes the same attack, altering only packet timing and sizes while preserving protocol compliance and attack functionality. The modified traffic is then classified as benign, allowing the attack to evade detection (bottom path). DEPT aims to detect and reverse these manipulations, restoring the NIDS's ability to identify the attack.}
    \label{fig:intro-evasion}
\end{figure}

Evaluating adversarial robustness requires accounting for the constraints of real network traffic.
In the image
domain, where adversarial examples were first studied, a perturbed pixel array can still be interpreted as an image; a perturbed feature vector of network traffic may not correspond to any realizable packet sequence, since packets must
keep their checksums, lengths, and protocol semantics
consistent~\cite{pierazzi2020intriguing,hashemi2019towards}. 
In most practical scenarios, an attacker operates in the \emph{black-box} setting: it can probe a deployed detector and observe
outcomes, but has access to neither the model's architecture nor its
gradients~\cite{papernot2017practical,apruzzese2023sok}. Much of the defensive literature
sets aside one or both of these constraints, so its robustness results may not reflect threats feasible under these conditions; \Cref{chap:related} examines this mismatch.

This work, together with a companion thesis, addresses this mismatch.
TANTGen~\cite{pissarra2026tantgen}
addresses the realism gap on the attack side: it is a Transformer-based generator that
manipulates attack traffic under black-box constraints, altering only packet timing
and packet size, and the traffic it emits is protocol-compliant by construction. DEPT provides
the defense side of the same effort. It is trained on, and evaluated against, the
realizable adversarial traffic that TANTGen produces, grounding the evaluation in packet-level attacks rather than feature perturbations whose realizability has not been established.

% #############################################################################
\section{Problem Statement}\label{sec:intro-problem}
This thesis addresses the following problem.
A stream of flows arrives at the
boundary of a protected network, where a \ac{NIDS} is deployed. Part of the traffic is
benign; part belongs to conventional attacks, which the deployed \ac{NIDS} is assumed to
detect; and part consists of attack traffic that has been adversarially manipulated to evade detection.
The defense must decide, for each flow and while the
flow is still in progress, whether it has been manipulated, and, when it has, counteract
the manipulation before the protected \ac{NIDS} processes the traffic.

Three constraints define an acceptable solution. First, \emph{streaming operation}: the
verdict must be produced packet by packet, as the flow unfolds, because waiting for a flow to complete may allow an attack to succeed before detection, while buffering entire flows can impose excessive memory and latency costs. 
Second, \emph{realistic attack assumptions}:
the defense must be trained and evaluated against realizable, protocol-compliant adversarial
traffic generated under the assumed black-box access constraints, rather than unconstrained feature-space perturbations.
Third, \emph{realizable
output}: a flow the defense repairs must remain valid network traffic, preserving protocol semantics as well as packet-level correctness; otherwise the defense has merely moved the
realizability problem from the attacker to the defender.

Two scope limitations complete the problem statement. DEPT complements the protected \ac{NIDS}, it does
not replace it: recognizing the underlying attacks remains the downstream system's task,
while DEPT detects and attempts to reverse adversarial manipulation of the traffic.
Since the defense is designed to operate inline, its processing overhead must also meet real-time latency requirements; evaluating the implemented system against those requirements is left to future work
(\Cref{sec:concl-future}).

% #############################################################################
\section{Work Objectives and Challenges}\label{sec:intro-objectives}
The problem statement is divided into four objectives, which structure both the design
presented in \Cref{chap:solution} and the research questions that guide the evaluation of
\Cref{chap:evaluation}:

\begin{itemize}
    \item \textbf{O1 (traffic representation).} Develop a representation of raw
    traffic that preserves sufficient packet-level information for model outputs to be
    mapped back to realizable traffic, and that architecturally different learners can
    consume at different temporal granularities, from individual packets to aggregated
    flow statistics.
    \item \textbf{O2 (heterogeneous detection).} Build a detection ensemble whose
    members belong to different architectural families, with the aim of reducing the transferability of adversarial perturbations between them, and combine their
    predictions into a single per-packet verdict.
    \item \textbf{O3 (input purification).} Learn to reconstruct a flagged flow:
    remove the adversarial manipulation and restore the flow's characteristics before manipulation,
    while preserving a valid, protocol-compliant packet sequence.
    \item \textbf{O4 (streaming-aware evaluation).} Evaluate the system under streaming conditions: on partial flows as they arrive, with metrics that measure how early
    a verdict becomes reliable, rather than end-of-flow accuracy alone.
\end{itemize}

Three challenges arise in this setting. The first is imbalance: the corpus used in this work is dominated
by high-rate floods, so a handful of traffic sources contributes the overwhelming
majority of flows. 
Training and evaluation must account for this imbalance to avoid under-representing other sources and obscuring poor performance on them on aggregate results. 
The second is the
limited observable signal: a manipulation confined to packet timing on a high-rate flood
may produce only small changes in the observed features, making manipulated and unmodified traffic difficult to distinguish; \Cref{chap:evaluation} quantifies this limitation.
The third is that the defense
shares the attacker's realizability constraint: purification edits traffic, and edited traffic must
remain realizable, which rules out treating flows as unconstrained feature vectors and informs the
representation choices of O1. All three challenges must also be addressed within the latency requirements of streaming operation.

% #############################################################################
\section{Rationale and Contributions}\label{sec:intro-contributions}
The design of DEPT rests on two decisions, motivated briefly here and discussed in more detail in
\Cref{chap:back,chap:related}.
The first
concerns detection. One strategy available to a black-box attacker exploits transferability:
perturbations crafted against a substitute model may also evade the target model~\cite{papernot2017practical}.
DEPT's detector therefore stacks two learners from different architectural families, a
\ac{RF}, which is not directly amenable to gradient-based optimization, and a \ac{TCN}, which processes packet-level feature sequences.
A small attention-based meta-learner combines their predictions at every packet~\cite{wolpert1992stacked}. 
This architectural diversity aims to reduce transferability: a perturbation that changes one learner's prediction may not change the other's. 

The second decision concerns what to do with a flagged flow. Dropping it would risk denying service to legitimate clients whose traffic is incorrectly flagged, so DEPT instead attempts to remove the manipulation through input purification, an approach explored in works such as
Defense-GAN~\cite{samangouei2018defense}. The purifier is a causal Transformer that
predicts a residual correction over the packet sequence, rather than a \ac{GAN} or a
diffusion model: self-attention captures dependencies between packets, causality allows each correction to depend only on the current and preceding packets, and a residual head allows the input to remain unchanged when the predicted correction is zero.

Building on this rationale, the thesis makes five contributions:

\begin{itemize}
    \item A packet-to-tensor pipeline that turns raw captures into masked,
    bit-level packet representations and aggregate flow statistics, so
    that architecturally different learners may process the same traffic at different
    temporal granularities (\Cref{sec:sol-pipeline}).
    \item A heterogeneous detection ensemble comprising an \ac{RF} and a \ac{TCN}, whose predictions are combined per packet by an attention-based meta-learner, together with a streaming comparison
    that evaluates every model by how early, and not only how well, it separates
    manipulated from unmodified flows (\Cref{sec:sol-ensemble,sec:eval-meta}).
    \item A causal purification Transformer coupled to an explicit realization step:
    flagged flows are corrected in feature space, the corrections are applied to packet timing and size, and the resulting packets are checked for parsing errors with a standard protocol dissector (\texttt{tshark})
    (\Cref{sec:sol-reconstructor,sec:eval-purification}).
    \item An end-to-end evaluation that measures DEPT as a defense placed in
    front of a protected NIDS: unmodified, TANTGen-perturbed, and DEPT-defended traffic are
    each evaluated by the independent NIDS targeted by the attack, reporting changes in the attack success and detection rates (\Cref{sec:eval-e2e}).
    \item An adversarial evaluation corpus of ten traffic sources drawn from three public
    datasets
    (Kitsune, CIC-IDS2017, and CIC-DDoS-2019~\cite{mirsky2018kitsune,cicids2017,sharafaldin2019cicddos}),
    each manipulated by TANTGen in three perturbation modes, with a per-source analysis
    of where the defense is effective, where it fails, and the factors underlying these outcomes (\Cref{sec:sol-pipeline} and
    \Cref{chap:evaluation}).
\end{itemize}

% #############################################################################
\section{Organization of the Document}\label{sec:intro-organization}
The remainder of this document is organized as follows. \Cref{chap:back} presents the
background: learning-based intrusion detection and traffic representation, adversarial
machine learning and the realizability constraints that distinguish it in the network
domain, and the ensemble concepts behind DEPT's detector. \Cref{chap:related} surveys
related work on the attack and the defense sides, and states the gap DEPT fills. \Cref{chap:solution} presents DEPT:
the threat model, a system overview, the corpus and packet-to-tensor pipeline, the
detection ensemble, and the purification and realization stages. \Cref{chap:evaluation}
evaluates the proposed system. \Cref{chap:conclusion}
closes with the answers to the research questions, the system's limitations, and future
work.
